\documentclass[journal,12pt,twocolumn]{IEEEtran}

\usepackage{setspace}
\usepackage{gensymb}
\usepackage{xcolor}
\usepackage[cmex10]{amsmath}
\usepackage{mathtools}
\usepackage{hyperref}
\usepackage{amsthm}
\usepackage{mathrsfs}
\usepackage{txfonts}
\usepackage{stfloats}
\usepackage{cite}
\usepackage{cases}
\usepackage{subfig}
\usepackage{longtable}
\usepackage{multirow}
\usepackage{enumitem}
\usepackage{listings}
\usepackage[section]{placeins}
\usepackage{needspace}
\usepackage{float}
\usepackage{graphicx}

\let\vec\mathbf
\DeclareMathOperator*{\Res}{Res}

%creating image environment
\newcommand{\image}[1]{\IfFileExists{#1}{\includegraphics[width=\columnwidth]{#1}}{\fbox{\parbox[c][0.85in][c]{0.88\columnwidth}{\centering Figure file not included in this source package.}}}}




\renewcommand\thesection{\arabic{section}}
\renewcommand\thesubsection{\thesection.\arabic{subsection}}
\renewcommand\thesubsubsection{\thesubsection.\arabic{subsubsection}}

\hyphenation{op-tical net-works semi-conduc-tor}

\lstset{
    language=Python,
    frame=single,
    breaklines=true,
    columns=fullflexible
}

\theoremstyle{definition}
\newtheorem{theorem}{Theorem}[section]
\newtheorem{problem}{Problem}
\newtheorem{proposition}{Proposition}[section]
\newtheorem{lemma}{Lemma}[section]
\newtheorem{corollary}[theorem]{Corollary}
\newtheorem{example}{Example}[section]
\newtheorem{definition}{Definition}[section]

\newcommand{\BEQA}{\begin{eqnarray}}
\newcommand{\EEQA}{\end{eqnarray}}
\newcommand{\define}{\stackrel{\triangle}{=}}
\newcommand{\myvec}[1]{\ensuremath{\begin{pmatrix}#1\end{pmatrix}}}
\newcommand{\mydet}[1]{\ensuremath{\begin{vmatrix}#1\end{vmatrix}}}

\providecommand{\nCr}[2]{\,^{#1}C_{#2}}
\providecommand{\nPr}[2]{\,^{#1}P_{#2}}
\providecommand{\mbf}{\mathbf}
\providecommand{\pr}[1]{\ensuremath{\Pr\left(#1\right)}}
\providecommand{\qfunc}[1]{\ensuremath{Q\left(#1\right)}}
\providecommand{\sbrak}[1]{\ensuremath{{}\left[#1\right]}}
\providecommand{\lsbrak}[1]{\ensuremath{{}\left[#1\right.}}
\providecommand{\rsbrak}[1]{\ensuremath{{}\left.#1\right]}}
\providecommand{\brak}[1]{\ensuremath{\left(#1\right)}}
\providecommand{\lbrak}[1]{\ensuremath{\left(#1\right.}}
\providecommand{\rbrak}[1]{\ensuremath{\left.#1\right)}}
\providecommand{\cbrak}[1]{\ensuremath{\left\{#1\right\}}}
\providecommand{\lcbrak}[1]{\ensuremath{\left\{#1\right.}}
\providecommand{\rcbrak}[1]{\ensuremath{\left.#1\right\}}}
\providecommand{\rect}[1]{\text{rect}\ensuremath{\left(#1\right)}}
\providecommand{\sinc}[1]{\text{sinc}\ensuremath{\left(#1\right)}}
\theoremstyle{remark}
\newtheorem{rem}{Remark}
\newcommand{\sgn}{\mathop{\mathrm{sgn}}}
\providecommand{\abs}[1]{\left\vert#1\right\vert}
\providecommand{\res}[1]{\Res\displaylimits_{#1}}
\providecommand{\norm}[1]{\lVert#1\rVert}
\providecommand{\mtx}[1]{\mathbf{#1}}
\providecommand{\mean}[1]{E\left[ #1 \right]}
\providecommand{\fourier}{\overset{\mathcal{F}}{\rightleftharpoons}}
\providecommand{\ztrans}{\overset{\mathcal{Z}}{\rightleftharpoons}}
\providecommand{\system}[1]{\overset{\mathcal{#1}}{\longleftrightarrow}}

\newcommand{\solution}{\noindent \textbf{Solution: }}
\providecommand{\dec}[2]{\ensuremath{\overset{#1}{\underset{#2}{\gtrless}}}}

\numberwithin{equation}{section}
\makeatletter
\@addtoreset{figure}{problem}
\makeatother

\begin{document}

\title{Charger}
\author{G V V Sharma$^{*}$
\thanks{$^{*}$The author is with the Department of Electrical Engineering,
Indian Institute of Technology Hyderabad, 502285, India.
E-mail: gadepall@iith.ac.in. All content in the manuscript is released
under GNU GPL. Free to use for anything.}}
\maketitle

\begin{abstract}
In this report, we use a full-wave-rectified signal to review periodic signals, Fourier series, Fourier transforms, and filtering. We then describe the main stages of a regulated DC supply and present our measured USB output values.
\end{abstract}

\section{Periodic Function}
\begin{enumerate}[label=\thesection.\arabic*.,ref=\thesection.\theenumi]
\item Consider the full-wave-rectified sinusoidal signal

\begin{align}
    x(t) &= A_0\abs{\sin\brak{2\pi f_0 t}},
    \label{eq:10-orig-diff-def}
\end{align}
Here, $A_0$ is the peak amplitude and $f_0$ is the frequency of the original sinusoid. Taking the absolute value turns both half-cycles positive, so the signal repeats twice during each cycle of the original sinusoid.

\begin{enumerate}[label=\thesection.\arabic*,
                  ref=\thesection.\theenumi]
    \item Plot $x(t)$.
    \item Show that $x(t)$ is periodic and find its period. In particular, distinguish the period of $x(t)$ from the period of the underlying sinusoid.
\end{enumerate}
\end{enumerate}

\FloatBarrier
\section{Fourier Series}
For our numerical calculations, we use $A_0=12\,\mathrm{V}$ and $f_0=50\,\mathrm{Hz}$. The rectified signal therefore has a fundamental frequency of $2f_0=100\,\mathrm{Hz}$.
\begin{enumerate}[label=\thesection.\arabic*,
                  ref=\thesection.\theenumi]
    \item If
    \begin{align}
        x(t) &= \sum_{k=-\infty}^{\infty}c_ke^{j2\pi kf_0t}
        \label{eq:one-Z-complex}
    \end{align}
    show that
    \begin{align}
        c_k &= f_0\int_{-\frac{1}{2f_0}}^{\frac{1}{2f_0}}
        x(t)e^{-j2\pi kf_0t}\,dt
        \label{eq:one-Z}
    \end{align}

    \item Find $c_k$ for \eqref{eq:10-orig-diff-def}. In particular, account for the fact that the signal has period $1/(2f_0)$, although the complex series is written using the underlying frequency $f_0$.

    \item Verify \eqref{eq:10-orig-diff-def} using Python.

    \item Show that
    \begin{align}
        x(t) &= \frac{a_0}{2}+\sum_{k=1}^{\infty}
        \left(a_k\cos(2\pi kf_0t)+b_k\sin(2\pi kf_0t)\right)
        \label{eq:one-Z-real}
    \end{align}
    and obtain the coefficients
    \begin{align*}
        a_0 &= 2f_0\int_{-\frac{1}{2f_0}}^{\frac{1}{2f_0}}x(t)\,dt,\\
        a_k &= 2f_0\int_{-\frac{1}{2f_0}}^{\frac{1}{2f_0}}
        x(t)\cos(2\pi kf_0t)\,dt,\\
        b_k &= 2f_0\int_{-\frac{1}{2f_0}}^{\frac{1}{2f_0}}
        x(t)\sin(2\pi kf_0t)\,dt,\quad k\\geq 1.
    \end{align*}

    \item Find $a_k$ and $b_k$ for \eqref{eq:10-orig-diff-def}. Use the same coefficient convention as in the real Fourier series above.

    \item Verify \eqref{eq:one-Z-real} using Python.
\end{enumerate}

For this signal, we can check the coefficients using these results:
$b_k=0$ for every $k$, $a_0/2=2A_0/\pi$, and
$a_{2n}=-4A_0/[\pi(4n^2-1)]$ for $n\geq1$, while $a_k=0$
for odd $k$. Equivalently, $c_0=2A_0/\pi$,
$c_{2n}=c_{-2n}=-2A_0/[\pi(4n^2-1)]$, and all odd-indexed
complex coefficients vanish.

\FloatBarrier
\section{Fourier Transform}
\begin{enumerate}[label=\thesection.\arabic*,
                  ref=\thesection.\theenumi]
    \item The Dirac delta is a generalized function characterized by
    \begin{align}
        \delta(t)&=0,\quad t\neq 0,\\
        \int_{-\infty}^{\infty}\delta(t)\,dt&=1,\\
        \int_{-\infty}^{\infty}g(t)\delta(t-t_0)\,dt&=g(t_0).
    \end{align}

    \item The Fourier Transform of $g(t)$ is
    \begin{align}
        G(f)=\int_{-\infty}^{\infty}g(t)e^{-j2\pi ft}\,dt.
    \end{align}

    \item Show that
    \begin{align}
        g(t-t_0)&\system{F}G(f)e^{-j2\pi ft_0}.
    \end{align}

    \item Show that
    \begin{align}
        G(t)&\system{F}g(-f).
    \end{align}

    \item Determine the Fourier transform of $\delta(t)$, treating it as a generalized function (distribution).

    \item Determine the Fourier transform of $e^{-j2\pi f_0t}$.

    \item Determine the Fourier transform of $\cos(2\pi f_0t)$.

    \item Find the Fourier transform of the periodic signal $x(t)$ and plot its line spectrum.
    Use the distributional Fourier-transform convention and verify the result using Python.

    \item Using the normalized definitions
    $\rect{t}=1$ for $|t|<1/2$ and 
    $\sinc{t}=\sin(\pi t)/(\pi t)$, show that
    \begin{align}
        \rect{t}\system{F}\sinc{f}.
    \end{align}
    Verify using Python.

    \item Determine the Fourier transform of $\sinc{t}$ using the same normalized convention.
    Verify using Python.
\end{enumerate}

\FloatBarrier
\section{Filter}
\begin{enumerate}[label=\thesection.\arabic*,
                  ref=\thesection.\theenumi]
    \item Determine the required DC gain and the harmonic-rejection condition on $H(f)$ so that the steady-state output is a constant $5\,\mathrm{V}$. State any assumptions about the filter passband.
    \item For an ideal low-pass model with cutoff $f_c$ satisfying $0<f_c<100\,\mathrm{Hz}$, find the corresponding impulse response $h(t)$.
    \item Verify your result using convolution.
\end{enumerate}

The input's average value is $X_{\mathrm{DC}}=2A_0/\pi$.
For $A_0=12\,\mathrm{V}$, the required DC gain is therefore
$K=5/(24/\pi)=5\pi/24\approx0.6545$. An ideal low-pass response
with $H(f)=K$ for $|f|<f_c$ and $H(f)=0$ for $|f|>f_c$, where
$0<f_c<100\,\mathrm{Hz}$, passes the DC component and rejects the
first and all higher harmonics. With the normalized sinc definition,
its impulse response is $h(t)=2Kf_c\sinc{(2f_ct)}$.

\FloatBarrier
\section{Filter Design}
\begin{enumerate}[label=\thesection.\arabic*,
                  ref=\thesection.\theenumi]
    \item Design a Butterworth approximation to the specified low-pass response $H(f)$. State the filter order, cutoff frequency, and any design tolerance used.
    \item Design a Chebyshev approximation to the specified low-pass response $H(f)$. State the filter order, passband ripple, and cutoff frequency.
    \item Design a circuit for your Butterworth filter.
    \item Design a circuit for your Chebyshev filter.
\end{enumerate}


\FloatBarrier
\section{Circuit Input and Transformer}
Our calculations use a $220\,\mathrm{V}$ RMS, $50\,\mathrm{Hz}$
primary supply and the measured $10.53\,\mathrm{V}$ RMS secondary.
The transformer steps down the voltage and provides galvanic isolation.

The modified charger circuit is shown in Fig.~\ref{fig:modified-charger-circuit}.

\begin{figure}[H]
    \centering
    \image{Project/images/modifiedCircuitSchematic.jpeg}
    \caption{Modified circuit schematic for the charger.}
    \label{fig:modified-charger-circuit}
\end{figure}

\begin{table}[H]
\centering
\caption{Input and transformer values}
\label{tab:transformer-values}
\small
\begin{tabular}{|p{0.55\columnwidth}|c|}
\hline
\textbf{Quantity} & \textbf{Value}\\ \hline
Primary RMS voltage & $220\,\mathrm{V}$\\
Mains frequency & $50\,\mathrm{Hz}$\\
Primary peak voltage (calculated) & $311.1\,\mathrm{V}$\\
Measured secondary RMS voltage & $10.53\,\mathrm{V}$\\
Secondary peak voltage (calculated) & $14.89\,\mathrm{V}$\\
Approximate turns ratio & $20.9:1$\\
\hline
\end{tabular}
\end{table}

We calculate peak voltage as $\sqrt{2}$ times RMS voltage. The ideal
transformer voltage ratio also gives the turns ratio.

The separate centre-tapped example uses a $12$--$0$--$12\,\mathrm{V}$
transformer, a $12\,\mathrm{V}$ RMS winding output, and an
approximately $18\,\mathrm{V}$ peak. These values do not apply to the
$10.53\,\mathrm{V}$ RMS secondary used above.

\FloatBarrier
\section{Full-Wave Bridge Rectifier}
In our four-diode bridge, two diodes conduct during each half-cycle,
keeping current through the load in one direction and doubling the
ripple frequency. Using a $0.7\,\mathrm{V}$ drop per diode, we
calculate a $13.49\,\mathrm{V}$ peak, $8.59\,\mathrm{V}$ average,
and $9.54\,\mathrm{V}$ RMS. These are first-order estimates; diode
conduction and source impedance affect the actual waveform.

\begin{table}[H]
\centering
\caption{Bridge rectifier values}
\label{tab:bridge-values}
\small
\begin{tabular}{|p{0.55\columnwidth}|c|}
\hline
\textbf{Quantity} & \textbf{Value}\\ \hline
Measured secondary RMS voltage & $10.53\,\mathrm{V}$\\
Secondary peak voltage & $14.89\,\mathrm{V}$\\
Assumed forward drop per diode & $0.7\,\mathrm{V}$\\
Bridge peak voltage & $13.49\,\mathrm{V}$\\
Rectified average voltage & $8.59\,\mathrm{V}$\\
Rectified RMS voltage & $9.54\,\mathrm{V}$\\
Input frequency & $50\,\mathrm{Hz}$\\
Full-wave ripple frequency & $100\,\mathrm{Hz}$\\
Ripple period & $10\,\mathrm{ms}$\\
\hline
\end{tabular}
\end{table}

Two rectifier diode options are the 1N4007 diode
($1\,\mathrm{A}$ average forward-current class and
$1000\,\mathrm{V}$ repetitive peak reverse-voltage rating) and
the 1N5408 diode ($3\,\mathrm{A}$ average forward-current class).

\FloatBarrier
\section{Smoothing Capacitor}
The $2200\,\mu\mathrm{F}$ smoothing capacitor charges near each
rectified peak and supplies the load between peaks, reducing output
ripple.

For small ripple and nearly constant load current, we estimate
peak-to-peak ripple as
\[
V_{r(pp)}\approx\frac{I_L}{f_rC}.
\]
At $f_r=100\,\mathrm{Hz}$ and $C=2200\,\mu\mathrm{F}$:
\[
V_{r(pp)}\approx4.545 I_L\ \mathrm{V},
\]
where $I_L$ is in amperes. For a $13.49\,\mathrm{V}$ peak, the
approximate average is
\[
V_{\mathrm{DC}}\approx13.49-2.273I_L\ \mathrm{V}.
\]
These are calculated estimates, not measured values.

The separate centre-tapped example specifies a $100\,\mu\mathrm{F}$
filter capacitor and an approximately $18\,\mathrm{V}$ DC level.
This differs from the $2200\,\mu\mathrm{F}$ calculation above.

\FloatBarrier
\Needspace{8\baselineskip}
\section{Regulator and Buck Converter}
We compare two 5 V regulator options: the 7805 linear regulator and
the LM2596 step-down converter.

\subsection{7805 Linear Regulator}
The 7805 provides a fixed $5\,\mathrm{V}$ output. TI lists about
$2\,\mathrm{V}$ dropout at a $1\,\mathrm{A}$ test condition; the
actual value depends on the device, load, and temperature.

Its approximate power dissipation is
\[
P_D\approx(V_{\mathrm{IN}}-V_{\mathrm{OUT}})I_{\mathrm{OUT}}.
\]
The 7805 family has current limiting and thermal protection. Its
approximately $2.4\,\mathrm{A}$ current-limit figure is not a
continuous rating; operation above $1.5\,\mathrm{A}$ depends on the
device and cooling.

\subsection{LM2596 Buck Converter}
The LM2596 family specifies a $4.5$--$40\,\mathrm{V}$ input range,
up to $3\,\mathrm{A}$ output under stated conditions, and a
$150\,\mathrm{kHz}$ oscillator. A module's usable current also
depends on its components, board, and cooling.

For an ideal buck converter, the duty ratio is
\[
D=\frac{V_{\mathrm{OUT}}}{V_{\mathrm{IN}}}.
\]
A $12\,\mathrm{V}$ to $5\,\mathrm{V}$ conversion gives an ideal
duty ratio of about $41.7\%$. This is a reference calculation, not
our measured USB operating point.

\begin{table}[H]
\centering
\caption{Regulator comparison}
\label{tab:regulator-comparison}
\small
\begin{tabular}{|p{0.24\columnwidth}|p{0.30\columnwidth}|p{0.30\columnwidth}|}
\hline
\textbf{Feature} & \textbf{7805} & \textbf{LM2596}\\ \hline
Method & Linear & Switching buck\\
Output & Fixed $5\,\mathrm{V}$ & Fixed or adjustable\\
Main consideration & Headroom and heat & Ripple, layout, and heat\\
Key specifications & About $2\,\mathrm{V}$ dropout at stated $1\,\mathrm{A}$ test & $4.5$--$40\,\mathrm{V}$ input; up to $3\,\mathrm{A}$; $150\,\mathrm{kHz}$\\
\hline
\end{tabular}
\end{table}

\FloatBarrier
\section{USB Output Port}
The USB port provides power through VBUS and return through GND.
D+ and D- carry USB 2.0 data and can identify legacy charging
sources. Type-C advertises current through its CC pins; USB-PD also
requires a compatible controller and negotiation.

\begin{table}[H]
\centering
\caption{USB charging current values}
\label{tab:usb-current-values}
\small
\begin{tabular}{|p{0.48\columnwidth}|c|}
\hline
\textbf{Port or mode} & \textbf{Current value}\\ \hline
USB 2.0 SDP & Up to $500\,\mathrm{mA}$*\\
USB 3.x SDP & Up to $900\,\mathrm{mA}$*\\
BC1.2 CDP & Up to $1.5\,\mathrm{A}$\\
BC1.2 DCP & Up to $1.5\,\mathrm{A}$\\
USB Type-C current advertisement & $1.5\,\mathrm{A}$ or $3\,\mathrm{A}$ at $5\,\mathrm{V}$\\
\hline
\end{tabular}
\end{table}
\noindent\textit{*These are nominal maximum current values under the applicable USB specification and configuration/enumeration conditions; they are not guaranteed by the connector alone.}

The table lists legacy proprietary D+/D- signatures. They are
implementation-specific, not universal USB rules.

\begin{table}[H]
\centering
\caption{Example legacy D+/D- charging signatures}
\label{tab:legacy-usb-values}
\small
\begin{tabular}{|c|c|c|}
\hline
\textbf{D+} & \textbf{D-} & \textbf{Reported current}\\ \hline
$2.0\,\mathrm{V}$ & $2.0\,\mathrm{V}$ & Approx. $0.5\,\mathrm{A}$\\
$2.0\,\mathrm{V}$ & $2.7\,\mathrm{V}$ & Approx. $1.0\,\mathrm{A}$\\
$2.7\,\mathrm{V}$ & $2.0\,\mathrm{V}$ & Approx. $2.1\,\mathrm{A}$\\
$2.7\,\mathrm{V}$ & $2.7\,\mathrm{V}$ & Approx. $2.4\,\mathrm{A}$\\
\hline
\end{tabular}
\end{table}

\FloatBarrier
\Needspace{20\baselineskip}
\section{Measured USB Output}
We measured an average USB output of $4.94\,\mathrm{V}$ and
$0.60\,\mathrm{A}$ using a voltmeter and ammeter. The instrument
also displayed $16\,\mathrm{mAh}$, which we list separately because
mAh measures charge, not current.

\begin{table}[H]
\centering
\caption{Measured USB output values}
\label{tab:measured-usb-output}
\small
\begin{tabular}{|p{0.55\columnwidth}|c|}
\hline
\textbf{Parameter} & \textbf{Measured value}\\ \hline
Average output voltage & $4.94\,\mathrm{V}$\\
Average output current & $0.60\,\mathrm{A}$\\
Power output & $2.96\,\mathrm{W}$\\
\hline
\end{tabular}
\end{table}

\end{document}


